\section{Related work and limitations}
\label{sec:related}
Batch-invariant kernels address numerical variation in parallel reductions~\cite{thinkingmachines,layercast}. LLM-42 combines verified speculation with fixed-shape execution and rollback~\cite{llm42}; MarginGate studies margin-triggered verification~\cite{margingate}; Vosti specifies deterministic inference across kernels and stateful serving components~\cite{vosti}. These works address complementary control points. HeteroSpec focuses on a local graph representation that preserves a specified quantized reference while retaining batching outside attention.

Speculative decoding is not automatically faster after draft-model costs are included~\cite{illusion,lossless}. Our results separate the local QB/PQ layout comparison from the full speculative-versus-target-only comparison. This distinction is central to interpreting the RTX and B570 closure panels.

The study has important limits. The original CUDA results use a fixed three-family panel; the 72-context replay samples specific layers, phases and key extents. The later RTX closure uses four previously exposed inputs, one model pair, one device/date and k=1. Its new complete-service timer is not the same as the original resident-request timer. The B570 closure uses a different frontend, device, runtime, reference and context policy. Its compact records were reconstructed from GitHub, but the complete 2.27 GB RTX raw archive remains local and has not been externally authenticated in this review. No result establishes general equivalence, unseen-input behavior, a physical GPU bottleneck, energy efficiency or a new accelerator benefit.
